% GENERATED FILE. Edit canonical lessons, then run npm run generate:books.
% canonical-source-hash: fnv1a64:4a0c4de8734f2aa0
% canonical-lessons: ML-C509-tapanila, ML-C509-ardrata, ML-C509-usnam, ML-C509-kulir, ML-C509-cutakuka

\chapter{Heat and Cold}
\label{ch:ml-heat-and-cold}
\hlchaptermodality{509}

\begin{chapteropening}
\textbf{By the end of this chapter:} \emph{I can say temperature, humidity, sultry heat, a chill and to get hot.}

Complete the last lesson of chapter 509: I can say temperature, humidity, sultry heat, a chill and to get hot.
\end{chapteropening}

\section[\texorpdfstring{\ml{താപനില} (tāpanila)}{താപനില (tāpanila)}]{\ml{താപനില} (tāpanila) — temperature}

\label{lesson:ML-C509-tapanila}



\begin{quote}
Before the new one: say the Malayalam for a raincoat, then the Malayalam for to dry out.
\end{quote}

\subsection*{You'll want to know: \ml{താപനില}}
\textbf{\ml{താപനില}} — \emph{tāpanila} — \textquotedblleft{}temperature\textquotedblright{}.

\begin{quote}
\ml{ഇന്ന്} \ml{താപനില} \ml{മുപ്പതിന്} \ml{മുകളിൽ} \ml{ആണ്}.
\end{quote}

\begin{quote}
\emph{innŭ tāpanila muppatinŭ mukaḷil āṇŭ.}
\end{quote}

\begin{quote}
Today's temperature is above thirty.
\end{quote}

\ml{താപം}, heat, plus \ml{നില}, a level: the level of heat.

\begin{cousinweb}
\textbf{\ml{താപനില}} (\emph{tāpanila}) is a Malayalam compound built from \ml{താപം}, heat, from Sanskrit \emph{tāpa} (Monier-Williams), and the native \ml{നില}, a standing or level (DEDR 3675).
\end{cousinweb}

\begin{grammarlens}[title={What you've built}]
One more word for getting around.
\end{grammarlens}

\subsection*{Your turn}
\begin{itemize}
\raggedright
  \item \emph{Say it:} \emph{tāpanila}
  \item \emph{Say it:} \emph{tāpanila}, once more
  \item \emph{Say it:} \emph{innŭ tāpanila muppatinŭ mukaḷil āṇŭ.} — Today's temperature is above thirty.
  \item \emph{Recall:} say the Malayalam for a raincoat, then the Malayalam for to dry out, then say \emph{tāpanila} again
\end{itemize}

\begin{tcolorbox}[breakable,colback=teal!4,colframe=teal!35!black,title={Before you move on}]
What does \ml{താപനില} mean? (\textquotedblleft{}Temperature\textquotedblright{}.) Say it once more.
\end{tcolorbox}
\section[\texorpdfstring{\ml{ആർദ്രത} (ārdrata)}{ആർദ്രത (ārdrata)}]{\ml{ആർദ്രത} (ārdrata) — humidity}

\label{lesson:ML-C509-ardrata}



\begin{quote}
Before the new one: say the Malayalam for to dry out, then the Malayalam for temperature.
\end{quote}

\subsection*{You'll want to know: \ml{ആർദ്രത}}
\textbf{\ml{ആർദ്രത}} — \emph{ārdrata} — \textquotedblleft{}humidity\textquotedblright{}.

\begin{quote}
\ml{തീരദേശത്ത്} \ml{ആർദ്രത} \ml{കൂടുതലാണ്}.
\end{quote}

\begin{quote}
\emph{tīradēśattŭ ārdrata kūṭutalāṇŭ.}
\end{quote}

\begin{quote}
Humidity is higher in the coastal region.
\end{quote}

A formal word from the weather report; people say the air is heavy and you sweat.

\begin{cousinweb}
\textbf{\ml{ആർദ്രത}} (\emph{ārdrata}) is built from Sanskrit \emph{ārdra}, \textquotedblleft{}wet, moist\textquotedblright{} (Monier-Williams), with \emph{-tā}, the ending of abstract nouns.
\end{cousinweb}

\begin{grammarlens}[title={What you've built}]
One more word for getting around.
\end{grammarlens}

\subsection*{Your turn}
\begin{itemize}
\raggedright
  \item \emph{Say it:} \emph{ārdrata}
  \item \emph{Say it:} \emph{ārdrata}, once more
  \item \emph{Say it:} \emph{tīradēśattŭ ārdrata kūṭutalāṇŭ.} — Humidity is higher in the coastal region.
  \item \emph{Recall:} say the Malayalam for to dry out, then the Malayalam for temperature, then say \emph{ārdrata} again
\end{itemize}

\begin{tcolorbox}[breakable,colback=teal!4,colframe=teal!35!black,title={Before you move on}]
What does \ml{ആർദ്രത} mean? (\textquotedblleft{}Humidity\textquotedblright{}.) Say it once more.
\end{tcolorbox}
\section[\texorpdfstring{\ml{ഉഷ്ണം} (uṣṇaṁ)}{ഉഷ്ണം (uṣṇaṁ)}]{\ml{ഉഷ്ണം} (uṣṇaṁ) — sultry heat}

\label{lesson:ML-C509-usnam}



\begin{quote}
Before the new one: say the Malayalam for temperature, then the Malayalam for humidity.
\end{quote}

\subsection*{You'll want to know: \ml{ഉഷ്ണം}}
\textbf{\ml{ഉഷ്ണം}} — \emph{uṣṇaṁ} — \textquotedblleft{}sultry heat\textquotedblright{}.

\begin{quote}
\ml{മഴയ്ക്ക്} \ml{മുൻപ്} \ml{വല്ലാത്ത} \ml{ഉഷ്ണമായിരുന്നു}.
\end{quote}

\begin{quote}
\emph{maḻaykkŭ munpŭ vallātta uṣṇamāyirunnu.}
\end{quote}

\begin{quote}
Before the rain the heat was oppressive.
\end{quote}

Damp, sticky heat; \ml{ചൂട്} is any heat.

\begin{cousinweb}
\textbf{\ml{ഉഷ്ണം}} (\emph{uṣṇaṁ}) is borrowed from Sanskrit \emph{uṣṇa}, \textquotedblleft{}hot, warm\textquotedblright{} (Monier-Williams).
\end{cousinweb}

\begin{grammarlens}[title={What you've built}]
One more word for getting around.
\end{grammarlens}

\subsection*{Your turn}
\begin{itemize}
\raggedright
  \item \emph{Say it:} \emph{uṣṇaṁ}
  \item \emph{Say it:} \emph{uṣṇaṁ}, once more
  \item \emph{Say it:} \emph{maḻaykkŭ munpŭ vallātta uṣṇamāyirunnu.} — Before the rain the heat was oppressive.
  \item \emph{Recall:} say the Malayalam for temperature, then the Malayalam for humidity, then say \emph{uṣṇaṁ} again
\end{itemize}

\begin{tcolorbox}[breakable,colback=teal!4,colframe=teal!35!black,title={Before you move on}]
What does \ml{ഉഷ്ണം} mean? (\textquotedblleft{}Sultry heat\textquotedblright{}.) Say it once more.
\end{tcolorbox}
\section[\texorpdfstring{\ml{കുളിർ} (kuḷir)}{കുളിർ (kuḷir)}]{\ml{കുളിർ} (kuḷir) — a chill; coolness}

\label{lesson:ML-C509-kulir}



\begin{quote}
Before the new one: say the Malayalam for humidity, then the Malayalam for sultry heat.
\end{quote}

\subsection*{You'll want to know: \ml{കുളിർ}}
\textbf{\ml{കുളിർ}} — \emph{kuḷir} — \textquotedblleft{}a chill; coolness\textquotedblright{}.

\begin{quote}
\ml{പുലർച്ചെ} \ml{മലയിൽ} \ml{നല്ല} \ml{കുളിർ} \ml{ആയിരുന്നു}.
\end{quote}

\begin{quote}
\emph{pularcce malayil nalla kuḷir āyirunnu.}
\end{quote}

\begin{quote}
At dawn it was very chilly in the hills.
\end{quote}

A pleasant coolness, as of a morning breeze.

\begin{cousinweb}
\textbf{\ml{കുളിർ}} (\emph{kuḷir}) is native Dravidian, \textquotedblleft{}coolness, chill\textquotedblright{} (DEDR 1834).
\end{cousinweb}

\begin{grammarlens}[title={What you've built}]
One more word for getting around.
\end{grammarlens}

\subsection*{Your turn}
\begin{itemize}
\raggedright
  \item \emph{Say it:} \emph{kuḷir}
  \item \emph{Say it:} \emph{kuḷir}, once more
  \item \emph{Say it:} \emph{pularcce malayil nalla kuḷir āyirunnu.} — At dawn it was very chilly in the hills.
  \item \emph{Recall:} say the Malayalam for humidity, then the Malayalam for sultry heat, then say \emph{kuḷir} again
\end{itemize}

\begin{tcolorbox}[breakable,colback=teal!4,colframe=teal!35!black,title={Before you move on}]
What does \ml{കുളിർ} mean? (\textquotedblleft{}A chill; coolness\textquotedblright{}.) Say it once more.
\end{tcolorbox}
\section[\texorpdfstring{\ml{ചൂടാകുക} (cūṭākuka)}{ചൂടാകുക (cūṭākuka)}]{\ml{ചൂടാകുക} (cūṭākuka) — to get hot}

\label{lesson:ML-C509-cutakuka}



\begin{quote}
Before the new one: say the Malayalam for sultry heat, then the Malayalam for a chill.
\end{quote}

\subsection*{You'll want to know: \ml{ചൂടാകുക}}
\textbf{\ml{ചൂടാകുക}} — \emph{cūṭākuka} — \textquotedblleft{}to get hot\textquotedblright{}.

\begin{quote}
\ml{ഉച്ചയായപ്പോൾ} \ml{കാറിന്റെ} \ml{ഉള്ള്} \ml{വല്ലാതെ} \ml{ചൂടായി}.
\end{quote}

\begin{quote}
\emph{uccayāyappōḷ kāṟinṟe uḷḷŭ vallāte cūṭāyi.}
\end{quote}

\begin{quote}
By noon the inside of the car got very hot.
\end{quote}

Past: \ml{ചൂടായി}, \emph{cūṭāyi}.

\ml{ചൂട്}, heat, plus \ml{ആകുക}, to become. A person can also \ml{ചൂടാകുക}: get angry.

\begin{cousinweb}
\textbf{\ml{ചൂടാകുക}} (\emph{cūṭākuka}) is a Malayalam verb built from the native \ml{ചൂട്}, heat (DEDR 2654), and \ml{ആകുക}, to become.
\end{cousinweb}

\begin{grammarlens}[title={What you've built}]
One more word for getting around.
\end{grammarlens}

\subsection*{Your turn}
\begin{itemize}
\raggedright
  \item \emph{Say it:} \emph{cūṭākuka}
  \item \emph{Say it:} \emph{cūṭākuka}, once more
  \item \emph{Say it:} \emph{uccayāyappōḷ kāṟinṟe uḷḷŭ vallāte cūṭāyi.} — By noon the inside of the car got very hot.
  \item \emph{Recall:} say the Malayalam for sultry heat, then the Malayalam for a chill, then say \emph{cūṭākuka} again
\end{itemize}

\begin{tcolorbox}[breakable,colback=teal!4,colframe=teal!35!black,title={Before you move on}]
What does \ml{ചൂടാകുക} mean? (\textquotedblleft{}To get hot\textquotedblright{}.) Say it once more.
\end{tcolorbox}
